\section{Smalltalk-80: a VM especificada em livro}
\label{sec:hist-smalltalk}

\paragraph{Imagem e máquina virtual.}
O \enquote{Blue Book}, \citetitle{livro:goldberg1983}
\parencite{livro:goldberg1983}, dedica a quarta parte à implementação. O
capítulo~26 separa dois componentes: a \emph{virtual image}, com todos
os objetos do sistema, e a \emph{virtual machine}, que dá dinâmica a
eles. A tarefa de quem implementa o Smalltalk-80 numa máquina nova é
escrever a máquina virtual; a imagem é carregada nela.

\paragraph{Bytecodes de pilha.}
O mesmo capítulo descreve o caminho do código: o compilador traduz cada
método para \enquote{sequences of eight-bit instructions called
bytecodes}, que são instruções para um interpretador orientado a pilha.
O exemplo do livro é o método \texttt{center} de \texttt{Rectangle}
(\texttt{\^{}origin + corner / 2}), que vira seis bytecodes: empilha a
primeira variável de instância, empilha a segunda, envia \texttt{+},
empilha a constante 2, envia \texttt{/} e retorna o topo da pilha. Abaixo
do interpretador fica a \emph{object memory}, que guarda os objetos; os
capítulos~27 a~30 especificam interpretador e memória em detalhe.

Compare com qualquer listagem \texttt{javap} da Parte~II: instruções de um
byte, operandos na pilha, chamadas de método como instrução
(\texttt{invokevirtual}) e um heap com coleta de lixo por baixo. A JVM
não copiou o Smalltalk, mas a forma é a mesma.

\paragraph{Inline caches.}
O Smalltalk também foi onde a chamada de método ficou rápida.
\textcite{paper:self-pics} contam a ideia, atribuída a
\textcite{paper:smalltalk-80-deutsch-schiffman}: o tipo do receptor num
ponto de chamada quase nunca muda, então o endereço do método encontrado
pode ser guardado no próprio ponto de chamada, e o método chamado só
confere se o tipo ainda é o mesmo. Esse
\termo{inline cache}{Cache do resultado da busca de método guardado no
próprio ponto de chamada, no código do chamador.} acerta 95\% das vezes em
código Smalltalk, número que o paper reúne de três trabalhos anteriores
(um deles o de Deutsch e Schiffman).

\paragraph{O que isso explica na JVM.}
\begin{itemize}
\item O formato da VM: bytecode compacto de pilha, interpretador e
  memória de objetos gerenciada. A seção~\ref{sec:javap-basico} é a
  versão Java do capítulo~26 do Blue Book.
\item As chamadas virtuais baratas. Quando a seção~\ref{sec:jit} fala de
  pontos de chamada monomórficos e megamórficos, está falando de inline
  caches e do que acontece quando eles não bastam, o problema que o Self
  atacou em seguida.
\end{itemize}
